\subsubsection{Was auf dem Lichtkegel
scheitert}\label{was-auf-dem-lichtkegel-scheitert}

Bei \codeword{det P=0} gilt weiterhin \codeword{rank J\_P=2}, aber
\codeword{C\_PJ\_P=0}. Der Zähler \codeword{J\_PC\_P} hat Rang zwei und
sein Quadrat ist null. Er ist somit kein Projektor; der massive Ausdruck
lässt sich nicht durch bloßes Einsetzen von P²=0 fortsetzen.

Es fehlt auch nicht nur diese besondere Formel. Für
\codeword{P=diag(1,0)} fixieren reelle und imaginäre Nullrotationen den
Impuls. Ein kovarianter Rückweg R müsste beide intertwinen. Subtraktion
dieser beiden Gleichungen erzwingt \codeword{(I\_3⊗E)R=0}, wobei
\codeword{E=[[0,1],[0,0]]}. Zugleich faktorisiert C\_P durch
\codeword{I\_3⊗E}; deshalb wäre \codeword{C\_PR=0}, im Widerspruch zu
\codeword{C\_PR=I\_2}. Das vollständige endliche Gleichungssystem hat
Koeffizientenrang zehn und augmentierten Rang elf. Dies beweist:
\textbf{Kein linearer Lorentz- kovarianter Rückweg auf den ganzen
Weylraum kann am Nullimpuls nur von P abhängen.} Ein gewöhnlicher, durch
Basiswahl fixierter Rückweg existiert natürlich, besitzt aber nicht
diese Kovarianz.

\subsection{4. Ein expliziter Rückweg mit zusätzlicher
Referenz}\label{ein-expliziter-ruxfcckweg-mit-zusuxe4tzlicher-referenz}

Ein zweiter Lorentzvektor Q erlaubt

\[\adjustbox{max width=\linewidth}{$\displaystyle 
 K_{P,Q}=C_PJ_Q=\sigma I_2+Q\operatorname{adj}P,
 \qquad\sigma=\operatorname{tr}(Q\operatorname{adj}P),
$}\]

\[\adjustbox{max width=\linewidth}{$\displaystyle 
 \det K_{P,Q}=2\sigma^2+\det(P)\det(Q),\qquad
 R_{P,Q}=J_Q K_{P,Q}^{-1}.
$}\]

Für ein zukünftiges Null-P und ein zukünftiges zeitartiges Q ist
σ\textgreater0, also ist der Rückweg endlich. Bei
\codeword{P=diag(1,0), Q=I} lautet er konkret

\[\adjustbox{max width=\linewidth}{$\displaystyle 
 R_{P,Q}=\begin{pmatrix}
 2&0\\0&0\\0&1/2\\1&0\\0&0\\0&1
 \end{pmatrix},\qquad C_PR_{P,Q}=I_2.
$}\]

Wenn P \textbf{und} Q transformiert werden, gilt die volle
Lorentz-Kovarianz. Q kann zum Beispiel einen Beobachter oder eine
Präparationsreferenz beschreiben. Diese zusätzliche Referenz ist aber
ein echter Eingang; ihre Existenz oder Verfügbarkeit wird nicht aus W
gewonnen.

Die lokale erste Ableitung C\_P darf außerdem nicht mit ihrem Rückweg
verwechselt werden. Die Division durch P² beziehungsweise durch
impulsabhängige Ausdrücke ist im Ortsraum im Allgemeinen nichtlokal. Auf
einer vorgegebenen massiven Massenschale ist P²=m² konstant; diese
Massenschale wäre wiederum eine zusätzliche Kinetikaussage.

\subsection{5. Früher Kinetiktest auf genau diesem
Auslesekanal}\label{fruxfcher-kinetiktest-auf-genau-diesem-auslesekanal}

Die Rang-zwei-Auslese liefert einen Spinor mit zwei komplexen
Komponenten. Die masselose Weylgleichung ist dadurch noch nicht erfüllt.
Für nullartiges P legt \codeword{adj(P) C\_PT=0} eine weitere
unabhängige lineare Bedingung fest. Der erlaubte T-Unterraum hat dann
Dimension fünf, sein Auslesebild Dimension eins. Das ist zunächst nur
eine algebraische Feldgleichungs- bedingung und keine vollständige
Einteilchen-Konstruktion.

Noch entscheidender ist der Test des tatsächlich vorgeschlagenen
Produkts
\(T^{\alpha\beta\dot\gamma}=F^{\alpha\beta}\bar\psi^{\dot\gamma}\). Für
freie masselose Konstituenten seien

\[\adjustbox{max width=\linewidth}{$\displaystyle 
 P_F=\lambda\bar\lambda^T,\quad
 P_\psi=\mu\bar\mu^T,\quad
 F=\lambda\lambda^T,\quad\bar\psi=\bar\mu.
$}\]

Beide Impulse sind nullartig. Der Prüfer behandelt alle acht
Spinorkomponenten als unabhängige Symbole und zeigt die Polynomidentität

\[\adjustbox{max width=\linewidth}{$\displaystyle 
 \boxed{C_{P_F+P_\psi}(F\otimes\bar\psi)=0.}
$}\]

Die reale physikalische Einschränkung folgt durch komplexe Konjugation
der gepunkteten Partner. Der Beweis ist zugleich die Leibnizregel: Der
Ableitungsterm auf F verschwindet durch dessen freie Feldstärke-
gleichung, der Term auf ψbar durch dessen freie Weylgleichung. Die
Identität gilt auch für nichtkollineare Impulse. Ein expliziter Fall
besitzt sogar den zeitartigen Gesamtimpuls \codeword{P\_F+P\_ψ=I} und
liefert trotzdem Auslese null.

Damit ist ein klarer positiver Anschlussauftrag übrig: Ein
\textbf{nichtnuller} physischer Spinor aus diesem Kanal braucht eine
berechnete massive, wechselwirkende oder außerhalb der Massenschale
liegende Quelle, eine andere Kinetik oder einen anderen ausdrücklich
festgelegten Feldadapter. Ein eingeschobener freier
Feldstärke-/Weyl-Vertrag würde an diesem Test scheitern. Das ist kein
Ausschluss des nativen Modenmodells, dessen Lorentzkinetik noch gar
nicht festgelegt ist.

\subsection{6. Ladung und die zwei verschiedenen
Komposita}\label{ladung-und-die-zwei-verschiedenen-komposita}

Im nativen U(1)-Vertrag besitzen \codeword{f,b} die Ladungen −1,−2 und
ihre Adjungierten +1,+2. Daher gilt unverändert

{\def\LTcaptype{none} % do not increment counter
\begin{longtable}[]{@{}
  >{\raggedright\arraybackslash}p{(\linewidth - 4\tabcolsep) * \real{0.3000}}
  >{\raggedleft\arraybackslash}p{(\linewidth - 4\tabcolsep) * \real{0.4000}}
  >{\raggedright\arraybackslash}p{(\linewidth - 4\tabcolsep) * \real{0.3000}}@{}}
\toprule\noalign{}
\begin{minipage}[b]{\linewidth}\raggedright
Operator
\end{minipage} & \begin{minipage}[b]{\linewidth}\raggedleft
N-Ladung
\end{minipage} & \begin{minipage}[b]{\linewidth}\raggedright
Bei wörtlicher Feldzuordnung
\end{minipage} \\
\midrule\noalign{}
\endhead
\bottomrule\noalign{}
\endlastfoot
D∼b f† & −1 & Wenn b in (1,0) und f† in (0,1/2), dann gemischthändig;
der neue Impulsadapter ist anwendbar \\
χ†∼b†ηf† & +3 & Bei adjungiertem b und f ist \codeword{(0,1)⊗(0,1/2)}
gleichhändig; der konjugierte alte Algebraadapter ist möglich \\
Impuls oder Ableitung & 0 & Ändert keine dieser Ladungen \\
\end{longtable}
}

\subsubsection{Die Zuordnung muss zusätzlich den ursprünglichen Vertex
bestehen}\label{die-zuordnung-muss-zusuxe4tzlich-den-urspruxfcnglichen-vertex-bestehen}

Die vorige Tabelle ist eine bedingte Fallunterscheidung und noch keine
gemeinsame Lorentz-Deutung des Hamiltonoperators. Fordert man
zusätzlich, dass f \textbf{wörtlich ein linkes Weylfeld} ist und
\codeword{b†ff} ohne Ableitung ein Lorentzskalar wird, dann muss das als
b† bezeichnete \textbf{Feld} selbst den Typ (1,0) tragen. Die
Spin-1-Darstellung ist zu ihrer algebraischen Dualdarstellung
äquivalent. Ihre skalare Paarung ist in der symmetrischen Basis

\[\adjustbox{max width=\linewidth}{$\displaystyle 
 \eta_3=\begin{pmatrix}0&0&1\\0&-2&0\\1&0&0\end{pmatrix};
 \qquad L^T\eta_3+\eta_3L=0.
$}\]

Alle drei Generatorgleichungen werden geprüft. Dagegen besitzt
\codeword{(0,1)⊗(1,0)=(1,1)} keinen Skalar; das entsprechende lineare
Invariantensystem hat Rang neun auf neun Koeffizienten.

Unter dieser \textbf{zusätzlichen wörtlichen Vertexbedingung} vertauscht
sich die Benennung der zwei Fälle:

\begin{itemize}
\tightlist
\item
  \codeword{b†} ist links, sein adjungiertes Feld \codeword{b} rechts;
\item
  \codeword{D=b f†} ist damit \textbf{gleichhändig rechts}. Der
  konjugierte alte C/R-Adapter liest einen rechten Spinor aus;
\item
  \codeword{χ†=b†ηf†} ist \textbf{gemischthändig} und wird durch den
  hier konstruierten Impulsadapter in einen linken Spinor ausgelesen.
\end{itemize}

Eine rechte Kompositquelle kann zur rechten Seite einer Weylgleichung
\codeword{sigma-bar·partial ψ = Quelle} passen. Diese Kinetik muss
allerdings erst auf derselben Konstruktion gezeigt werden. Umgekehrt
führt die oben angesetzte Wahl \codeword{b} selbst links, bei wörtlicher
Adjunktion, zu einem gemischten D, aber nicht zu einem skalaren nackten
\codeword{b†ff}-Vertex.

Diese zusätzliche Prüfung verhindert, das Feldwörterbuch isoliert an
einem Komposit zu reparieren und dabei den ursprünglichen Vertex zu
verlieren. Da Fock-Modenkoeffizienten nicht bereits vollständige lokale
Felder sind, ist damit keine konsistente Moden-Feld-Abbildung
ausgeschlossen. Sie muss ihre Komponenten und Adjungierten ausdrücklich
angeben.

Die innere Metrik η vertauscht keine Lorentzhändigkeit. Die beiden
Komposita haben gleiche Ladung \textbf{modulo vier}, aber verschiedene
U(1)-Ladungen, Fockwirkungen und gegebenenfalls Lorentztypen. Ein
neutraler Ableitungsfaktor macht sie nicht identisch. Der frühere
Z4-erhaltende zusätzliche Hamiltonkanal kann sie koppeln; das ist eine
explizite Dynamikerweiterung.

Eine unabhängige Weylkomponente passender Händigkeit und Ladung oder
eine konsistente ladungskonjugierte Feldabbildung bleibt eine
Alternative. Sie muss jedoch mit den nativen CAR, Ladungen und
W-Kontraktionen vereinbar konstruiert werden. Eine punktierte Komponente
durch reine Notation in eine unpunktierte umzubenennen ist kein solcher
Adapter.

\subsection{7. Konkrete Voraussetzungen für die globale
Quell-Komposition}\label{konkrete-voraussetzungen-fuxfcr-die-globale-quell-komposition}

Die Kompositionssuche erhält aus dieser Rechnung unabhängige, früh
prüfbare Anforderungen - keinen bereits hergeleiteten Raum:

\begin{enumerate}
\def\labelenumi{\arabic{enumi}.}
\tightlist
\item
  \textbf{Feldindexvertrag:} Den genauen Operator mit U(1)-Ladung und
  gepunkteten/ungepunkteten Indizes festlegen und zugleich den skalaren
  Vertex einschließlich Adjunktion prüfen. Bei gleicher Händigkeit
  genügt algebraisch der alte C/R-Adapter. Bei gemischter Händigkeit
  wird mindestens ein zusätzliches Lorentzvektorobjekt benötigt.
\item
  \textbf{Herkunft von P:} Falls P aus der Komposition stammen soll,
  muss die Quelle ein neutrales Vierervektor-/Ableitungsobjekt und
  dessen Lorentztransformationsgesetz liefern. Innere 64/60-Labels,
  Bosonzahl, Clock oder ein bloßer Graphabstand erfüllen diese
  Gleichungen nicht durch ihre Bezeichnung. Ein translationsartiger
  Ansatz benötigt darüber hinaus seine konkreten Generatoren.
\item
  \textbf{Gemeinsamer Operator:} Auf derselben Konstruktion die oben
  angegebene C\_P-Matrix anwenden und die Intertwining-Gleichung
  einschließlich Boosts prüfen. Am Lichtkegel dürfen Auslese und
  kovarianter Rückweg nicht gleichgesetzt werden; ein Rückweg verlangt
  die nachgewiesene zusätzliche Referenz.
\item
  \textbf{Nichtnulltest nach Kinetik:} Vor jeder Pol- oder
  Teilchenbehauptung das konkrete Matrixelement des Ableitungskomposits
  berechnen. Der freie masselose Produktansatz liefert hier exakt null.
  Eine erfolgreiche Alternative muss diesen Test sichtbar bestehen.
\item
  \textbf{Verbindung zum geladenen nativen Pol:} Referenzzustand,
  N-Ladung, Dispersions-/Massenschalenannahme und spektrales Gewicht
  müssen dieselben sein. Die innere endliche Auslese beweist diese
  Verbindung nicht automatisch.
\end{enumerate}

Die endliche Quell-Kompositions- und Ablaufprüfung kann parallel dazu
fortschreiten. Ein fehlendes P verbietet ihre innere Algebra nicht; es
verhindert nur, diesen Stand schon als vollzogenen gemischthändigen
Lorentz-Ausleseadapter auszugeben.

\subsection{Reproduktion und
Quellenumfang}\label{reproduktion-und-quellenumfang}

\begin{verbatim}
python3 -B replay.py
\end{verbatim}

Der eigenständige Lauf benötigt Python und SymPy. Er prüft vier
eingefrorene Quellenpins und \textbf{85 exakte Bedingungen je Modus}. Es
gibt keine numerischen Prüfbedingungen. Normal und \codeword{-OO} müssen
byte-identische JSON-Dateien liefern; Warnungen gelten als Fehler, und
keine Bedingung verwendet abschaltbare Assertions.

Die v1.6.8-Auslesematrizen werden direkt aus dem eingefrorenen JSON
eingelesen und mit der neuen Konstruktion verbunden. Die frühere gesamte
Feldsuite wird nicht nochmals ausgeführt. Die vollständigen Matrizen,
Rangzeugen und Quellenpins stehen in \codeword{results.normal.json}, die
aktuellen Programm- und Ergebnishashes im
\codeword{replay\_manifest.json}.

\begin{center}\rule{0.5\linewidth}{0.5pt}\end{center}

\section{Historischer Teil: vollständige Hauptfassung
v1.6.8}\label{historischer-teil-vollstuxe4ndige-hauptfassung-v1.6.8}

Die gesamte bisherige Hauptfassung folgt wortgetreu. Aktuell maßgeblich
sind die vorangestellten Kapitel v1.6.9. Sie vereinfachen den früheren
Z4-Versuch durch Nb, präzisieren die Clock-Komposition und erweitern den
Lorentzadapter. Alte Vorhaben oder Bedingungen gelten dadurch nicht als
neu bewiesen. Die frühere vollständige Reproduktion ist als
unverändertes Archiv eingefroren.

\section{TFPT / Universalraum: gemeinsame Schatten und ein
ursprünglicher kausaler
Prozess}\label{tfpt-universalraum-gemeinsame-schatten-und-ein-urspruxfcnglicher-kausaler-prozess}

Vollständige Fortschreibung v1.6.8, 15. September 2026.
Forschungsbericht und unabhängige Prüfung der nachgereichten
Forschungssynthese 1.0. Keine vollständige TOE.

\subsection{1. Ergebnis und
Leseschlüssel}\label{ergebnis-und-leseschluxfcssel}

Die drei parallel bearbeiteten Aufgaben liefern diesmal konkrete
mathematische Anschlüsse: einen kleinen ursprünglichen
Singulettmechanismus, die vollständige bedingte Schattenkarte des hellen
Raums und einen kausalen Eingriff unter dem unveränderten
Hamiltonoperator. Die nachgereichte Forschungssynthese verbindet diese
Aufgaben sinnvoll durch gemeinsame Operationsgeschichten. Ihre
mitgelieferten Rechnungen sind reproduziert; ein Satz in Abschnitt 8.4
benötigt jedoch eine Korrektur.

Die noch später eingegangenen Texte liefern außerdem eine frisch
reproduzierte fünfte Bosonstufennorm. Eine unabhängig zertifizierte
größere Variationskompression verbessert die Grundenergieobergrenze auf
-1.13847609Δ und unter dem beibehaltenen N=63-Satz die untere
Polenergiegrenze auf 0.01657709Δ. Dazu kommt ein neuer exakter
Händigkeitstest des Feldprojektors. Der Nachtrag „Neue Norm, korrigierte
Zustandszuordnung und Feldadapter`` enthält diese maßgeblichen
Ergänzungen; ältere Ritzkennzahlen verschiedener Zustände werden dort
ausdrücklich getrennt.

Die zentrale Vereinfachung ist nicht eine neue große Geometrie, sondern
eine sorgfältig ausgewählte kleine gemeinsame Ausführung. Ein Zustand,
seine messbaren Bilder und seine späteren Antworten müssen aus denselben
Operatoren berechnet werden. Zusätzliche Zustände, Kontrollen und
Referenzen werden dabei ausdrücklich mitgezählt.

\textbf{Exakt} bezeichnet algebraische Identitäten, Invarianzbeweise und
rational/modular abgesicherte Ränge im angegebenen Modell.
\textbf{Numerisch} bezeichnet unabhängige Gleitkomma-Rekonstruktionen
und den vollständigen endlichen Dynamikvergleich. \textbf{Bedingt}
heißt: Der Satz stimmt, benötigt aber ausdrücklich gewährte
Präparationen, Kontrollen oder Messungen. \textbf{Offen} bleibt deren
ursprüngliche Herleitung und die physikalische Fortsetzung.
Prüfprogramme sind keine formale Beweisassistenten-Verifikation.

Die vollständige bisherige Hauptfassung v1.6.7 bleibt im historischen
Teil wortgetreu erhalten. Die neuen Kapitel, die vollständigen drei
Teilberichte, ihr gemeinsames Dreizustands-Addendum und alle drei
Nachtragsprüfungen stehen davor. Das Quellenpaper bleibt als
unveränderte Eingabe im Prüfpaket, nicht als still korrigierte
Originaldatei. Aktuell maßgeblich sind diese Einordnung, die Korrekturen
und die neu dokumentierten Ressourcenbedingungen.

\subsection{2. Prüfung des nachgereichten
Forschungspapers}\label{pruxfcfung-des-nachgereichten-forschungspapers}

Geprüft wurden das 33-seitige PDF, sein vollständiger Markdowntext,
Prüfprotokoll, Quellenmanifest und eigene Ergebnisdatei sowie das
zusätzlich vorhandene Reproduktionsarchiv. Die vollständige Synthese
wurde gelesen; nicht jede der 74 historischen Quellen wurde neu
bewiesen. Die relevanten PDF-Seiten 16 bis 20 und 22 wurden zusätzlich
bildlich geprüft.

{\def\LTcaptype{none} % do not increment counter
\begin{longtable}[]{@{}
  >{\raggedright\arraybackslash}p{(\linewidth - 2\tabcolsep) * \real{0.5000}}
  >{\raggedright\arraybackslash}p{(\linewidth - 2\tabcolsep) * \real{0.5000}}@{}}
\toprule\noalign{}
\begin{minipage}[b]{\linewidth}\raggedright
Behauptung oder Baustein
\end{minipage} & \begin{minipage}[b]{\linewidth}\raggedright
Ergebnis der unabhängigen Prüfung
\end{minipage} \\
\midrule\noalign{}
\endhead
\bottomrule\noalign{}
\endlastfoot
74 eingefrorene Quellen & Alle Archivgrößen und SHA-256-Werte stimmen
mit dem Manifest überein \\
Markdown, PDF und TeX & Gelieferte bzw. archivierte Dateien stimmen mit
dem Prüfprotokoll überein \\
Neuester Stand S074 & Identisch mit der vollständigen hiesigen v1.6.7
einschließlich Prozess-/Schattennachtrag \\
57 eigene Prüfbedingungen & Aus dem isoliert entpackten Paket normal und
optimiert wiederholt; identische Ergebnisbytes \\
20 übernommene Kettenprüfungen & Ebenso wiederholt; identische
Ergebnisbytes \\
Grundzustand und frühere 822er-Suite & In dieser Runde nicht vollständig
erneut ausgeführt; keine entsprechende Behauptung \\
Wortkernrekonstruktion & Unter den angegebenen Positivitäts-, Relations-
und Beschränktheitsvoraussetzungen korrekt \\
Aussage zu der ersten Gramzeile in §8.4 & Zu korrigieren; siehe unten \\
Ein-Cartan-Referenz und Projektorgeometrie & Korrekte bedingte
Konstruktionen, kein nativer physikalischer Ursprung \\
\end{longtable}
}

